\documentclass[11pt]{article}
\usepackage[margin=1in]{geometry}
\usepackage{amsmath,amssymb,amsthm}
\usepackage{graphicx}
\usepackage[hidelinks]{hyperref}

\newtheorem{theorem}{Theorem}
\newtheorem{lemma}{Lemma}
\newtheorem{corollary}{Corollary}

\newcommand{\ACK}{\operatorname{ACK}}
\newcommand{\Ext}{\operatorname{Ext}}
\newcommand{\aco}{\operatorname{aco}}
\newcommand{\dist}{\operatorname{dist}}
\newcommand{\Lin}{\mathcal L}
\newcommand{\T}{\mathbb T}
\newcommand{\eps}{\varepsilon}

\title{A Uniform Subcase of \(\ell_\infty\)-Stability for Quasi-\(\ACK\) Structure}
\author{Run fa\_banach\_001, agent\_lane\_10}
\date{June 16, 2026}

\begin{document}
\maketitle

\section*{Source question}

The source is Geunsu Choi and Mingu Jung,
\emph{A generalized ACK structure and the denseness of norm attaining
operators}, arXiv:2204.01991.  In Remark 3.4, after proving that \(c_0(Z)\)
has quasi-\(\ACK\) structure whenever \(Z\) has quasi-\(\ACK\) structure, the
paper states that it is unknown whether quasi-\(\ACK\) structure is stable
under \(\ell_\infty\)-sum operations in general.  The same difficulty is
repeated after Theorem 4.6 for arbitrary families.

\begin{center}
\includegraphics[width=0.95\textwidth]{figures/open_problem_crop.png}
\end{center}

\section*{Definitions used in the packet}

We use the terminology of the source.  A quasi-\(\ACK\) structure on a Banach
space \(X\) consists of a \(1\)-norming set \(\Gamma\subseteq B_{X^*}\), a
function \(\rho:\Gamma\to[0,1)\), and, for each
\(e^*\in\Ext(B_{X^*})\), a subset \(\Gamma_{e^*}\subseteq \Gamma\) such that
\[
e^*\in \overline{\T\Gamma_{e^*}}^{w^*},\qquad
\sup_{x^*\in\Gamma_{e^*}}\rho(x^*)<1,
\]
and the local \(\ACK\)-type conditions (i)'--(vi)' of Definition 2.1 in the
source hold on every nonempty relatively weak-star open subset of
\(\Gamma_{e^*}\).

For such fixed data, call
\[
\Delta=\left(\bigcup_{e^*\in\Ext(B_{X^*})}\Gamma_{e^*}\right)\setminus\{0\}
\]
the active part of the quasi-\(\ACK\) witness.

\section*{Result}

\begin{theorem}\label{thm:main}
Let \(\{X_i\}_{i\in I}\) be a family of Banach spaces.  Suppose that each
\(X_i\) has quasi-\(\ACK\) data \((\Gamma_i,\rho_i,\Gamma_{i,e^*})\), and put
\[
\Delta_i=\left(\bigcup_{e^*\in\Ext(B_{X_i^*})}\Gamma_{i,e^*}\right)\setminus\{0\}.
\]
Assume the active parts are uniformly bounded away from \(1\):
\[
r:=\sup_{i\in I}\sup_{x^*\in\Delta_i}\rho_i(x^*)<1.
\]
Then
\[
Z=\left(\bigoplus_{i\in I}X_i\right)_{\ell_\infty}
\]
has \(\ACK_r\) structure.  In particular, \(Z\) has quasi-\(\ACK\) structure.
\end{theorem}

\section*{Idea of the proof}

The source's obstruction for \(\ell_\infty\)-sums is the lack of an explicit
description of the extreme points of the dual ball.  In the uniform subcase we
do not need to describe them.  Milman's converse to Krein-Milman shows that the
union of all active quasi-\(\ACK\) pieces is already norming.  If the
\(\rho\)-values on that union are uniformly below \(1\), the quasi-\(\ACK\)
local pieces can be used as ordinary \(\ACK_r\) pieces.

Once each summand is upgraded this way, the \(\ell_\infty\)-sum construction is
coordinate-local: choose a point of the weak-star open set in one coordinate,
apply the summand's \(\ACK_r\) data there, and extend the resulting vector and
operator by zero on all other coordinates.

\section*{Proof}

\begin{lemma}\label{lem:active}
Let \(X\) have quasi-\(\ACK\) data
\((\Gamma,\rho,\Gamma_{e^*})\), and let
\[
\Delta=\left(\bigcup_{e^*\in\Ext(B_{X^*})}\Gamma_{e^*}\right)\setminus\{0\}.
\]
If \(r=\sup_{\Delta}\rho<1\), then \(X\) has \(\ACK_r\) structure with
witnessing set \(\Delta\).
\end{lemma}

\begin{proof}
First \(\Delta\) is \(1\)-norming.  Indeed, each extreme point
\(e^*\in\Ext(B_{X^*})\) belongs to
\(\overline{\T\Gamma_{e^*}}^{w^*}\), hence to
\(\overline{\T\Delta}^{w^*}\).  By Krein-Milman,
\[
B_{X^*}=\overline{\operatorname{co}\Ext(B_{X^*})}^{w^*}
\subseteq \overline{\aco(\T\Delta)}^{w^*}.
\]
The reverse inclusion is immediate from \(\Delta\subseteq B_{X^*}\), so
\(B_{X^*}=\overline{\aco(\T\Delta)}^{w^*}\), which is equivalent to
\(\Delta\) being \(1\)-norming.

Let \(\eps>0\) and let \(U\) be a nonempty relatively \(w^*\)-open subset of
\(\Delta\).  Pick \(u^*\in U\).  By definition of \(\Delta\), there is an
extreme point \(e_0^*\in\Ext(B_{X^*})\) such that
\(u^*\in\Gamma_{e_0^*}\).  Then
\[
U_0=U\cap \Gamma_{e_0^*}
\]
is a nonempty relatively \(w^*\)-open subset of \(\Gamma_{e_0^*}\).  Apply the
quasi-\(\ACK\) condition to \(U_0\).  It gives a nonempty
\(V\subseteq U_0\), a point \(x_1^*\in V\), a vector \(e\in S_X\), and an
operator \(F\in\Lin(X)\) satisfying the local conditions (i)'--(vi)' relative
to the original set \(\Gamma\).  These same conditions hold a fortiori for the
smaller norming set \(\Delta\).  In condition (iv)' the bound is
\[
|x^*(F(e))|\leq \rho(x_1^*)\leq r
\]
for \(x^*\in\Delta\setminus V_1\).  Thus the \(\ACK_r\) conditions hold for
\(\Delta\).
\end{proof}

We now prove Theorem \ref{thm:main}.  By Lemma \ref{lem:active}, each \(X_i\)
has \(\ACK_r\) structure with the \(1\)-norming set \(\Delta_i\).
Let \(P_i:Z\to X_i\) and \(J_i:X_i\to Z\) denote the canonical projection and
injection.  Set
\[
\Delta=\bigcup_{i\in I}P_i^*(\Delta_i)\subseteq B_{Z^*}.
\]
This set is \(1\)-norming for \(Z\), because for \(z=(z_i)\in Z\),
\[
\sup_{\phi\in\Delta}|\phi(z)|
=\sup_{i\in I}\sup_{x^*\in\Delta_i}|x^*(z_i)|
=\sup_{i\in I}\|z_i\|
=\|z\|.
\]

Let \(U\) be a nonempty relatively \(w^*\)-open subset of \(\Delta\) and fix
\(P_{i_0}^*x_0^*\in U\), with \(x_0^*\in\Delta_{i_0}\).  Then
\[
U_{i_0}=\{x^*\in\Delta_{i_0}:P_{i_0}^*x^*\in U\}
\]
is a nonempty relatively \(w^*\)-open subset of \(\Delta_{i_0}\).  Apply
\(\ACK_r\) structure of \(X_{i_0}\) to \(U_{i_0}\) and \(\eps\), obtaining
\[
V_{i_0}\subseteq U_{i_0},\quad x_1^*\in V_{i_0},\quad
e_{i_0}\in S_{X_{i_0}},\quad F_{i_0}\in\Lin(X_{i_0}).
\]
Define
\[
V=P_{i_0}^*(V_{i_0}),\qquad z_1^*=P_{i_0}^*x_1^*,\qquad
e=J_{i_0}e_{i_0},\qquad F=J_{i_0}F_{i_0}P_{i_0}.
\]
Then \(\|F\|=\|F(e)\|=1\), \(z_1^*(F(e))=1\), and
\(F^*(z_1^*)=z_1^*\).

It remains only to check the three conditions involving all functionals in
\(\Delta\).  Let \(z^*=P_j^*y^*\in\Delta\).  If \(j=i_0\), the assertions are
exactly the \(\ACK_r\) assertions for \(X_{i_0}\).  If \(j\ne i_0\), then
\[
z^*(F(e))=0,\qquad F^*(z^*)=0.
\]
Hence condition (iv)' is automatic, and
\(\dist(F^*(z^*),\aco\{0,V\})=0\), which gives condition (v)'.  Condition
(vi)' concerns only the elements of \(V=P_{i_0}^*(V_{i_0})\), so it also
follows from the summand \(X_{i_0}\).  Therefore \(Z\) has \(\ACK_r\)
structure.

\section*{Limitations}

This is a partial result, not a solution of the full source problem.  The
general case remains open when the available quasi-\(\ACK\) witnesses are
genuinely non-uniform, i.e. when the active \(\rho\)-values cannot be bounded
away from \(1\).  The result also does not settle the source's separate
question about stability of property quasi-\(\beta\) under \(\ell_\infty\)
sums.

\section*{Novelty and literature check}

The cheap run indexes were searched for arXiv:2204.01991, generalized ACK,
quasi-\(\ACK\), property quasi-\(\beta\), norm attaining operators, and
\(\ell_\infty\)-sum stability.  No duplicate packet or attempt for this exact
target was found.  Web searches for the paper title, quasi-\(\ACK\) structure,
\(\ell_\infty\)-sums, and property quasi-\(\beta\) found the source paper but
no later resolution of the \(\ell_\infty\)-stability question in this bounded
pass.

\section*{Human-review recommendation}

This packet is a likely valid partial result.  Review should focus on Lemma
\ref{lem:active}, especially the use of Krein-Milman/Milman's converse to
show that the active union is norming, and on the coordinatewise verification
of condition (v)' in the \(\ell_\infty\)-sum.

\begin{thebibliography}{9}

\bibitem{ChoiJung}
G. Choi and M. Jung,
\emph{A generalized ACK structure and the denseness of norm attaining
operators}, arXiv:2204.01991.

\bibitem{Phelps}
R. R. Phelps,
\emph{Lectures on Choquet's Theorem}, second edition,
Lecture Notes in Mathematics 1757, Springer, 2001.

\end{thebibliography}

\end{document}
